\documentclass[11pt,a4paper]{article}
\usepackage{fontspec}
\usepackage{amsmath,amssymb,amsthm}
\usepackage{mathtools}
\usepackage{bm}
\usepackage{geometry}
\usepackage{enumitem}
\usepackage{xcolor}
\usepackage[pdfencoding=auto,psdextra]{hyperref}

\geometry{margin=2.5cm}
\setlength{\emergencystretch}{3em}

\newcommand{\E}{\mathbb{E}}
\newcommand{\R}{\mathbb{R}}
\newcommand{\KL}{\mathrm{KL}}
\newcommand{\N}{\mathcal{N}}
\newcommand{\Xspace}{\mathcal{X}}
\newcommand{\Za}{\mathcal{Z}_{A}}
\newcommand{\Zb}{\mathcal{Z}_{B}}
\newcommand{\Enc}{\mathcal{E}}
\newcommand{\Dec}{\mathcal{D}}
\newcommand{\Tmap}{\mathsf{T}}
\newcommand{\Tinv}{\mathsf{T}^{-1}}
\newcommand{\push}{\#}
\newcommand{\Jac}{\mathbf{J}}
\newcommand{\norm}[1]{\left\lVert #1 \right\rVert}

\theoremstyle{definition}
\newtheorem{definition}{Definition}
\newtheorem{problem}{Problem}
\newtheorem{remark}{Remark}

\title{\textbf{Cross-Latent Velocity-Field Transport for On-Policy Distillation}\\[0.3em]
\large A Problem Formalization}
\author{Jayce-Ping}
\date{}

\begin{document}
\maketitle

\section{Problem Setup}
\label{sec:setup}

\textbf{Two flow-matching generative models.}
Let $\Xspace\subseteq\R^{D}$ be a shared pixel (image) space, and let there be two
latent spaces induced by their respective variational autoencoders (VAEs),
\[
\Za\subseteq\R^{d_A},\qquad \Zb\subseteq\R^{d_B},
\]
equipped with encoder/decoder pairs $(\Enc_A,\Dec_A)$ and $(\Enc_B,\Dec_B)$, where
$\Enc_\bullet:\Xspace\to\mathcal{Z}_\bullet$ and $\Dec_\bullet:\mathcal{Z}_\bullet\to\Xspace$.
In general the two latent spaces are \emph{heterogeneous}:
\[
d_A\neq d_B,\quad \Za\not\cong\Zb,\quad \Enc_A\neq\Enc_B,
\]
i.e.\ the channel counts, spatial downsampling rates, and per-channel statistics may all differ.

On each latent space there is a (flow-matching / diffusion) generative model described by a
time-dependent velocity field:
\[
v_A:\Za\times[0,1]\to\R^{d_A},
\qquad
v_B:\Zb\times[0,1]\to\R^{d_B}.
\]
Each defines a probability-flow ODE that transports the prior at the noise end $t=1$ to the
latent marginal at the data end $t=0$:
\begin{equation}
\dot z_A(t) = v_A\big(z_A(t),t\big),
\qquad
\dot z_B(t) = v_B\big(z_B(t),t\big).
\label{eq:pf-ode}
\end{equation}
Denote the corresponding latent marginals by $p_A(\cdot,t)$ and $p_B(\cdot,t)$; the clean-end
($t=0$) marginals are $p_A^{0}=\Enc_A\push\, p_{\mathrm{data}}$ and
$p_B^{0}=\Enc_B\push\, p_{\mathrm{data}}$. After decoding, both generate the same pixel
distribution $p_{\mathrm{data}}$.

\textbf{Distillation direction.}
Fix one model as the \emph{teacher} (denoted $A$, with frozen velocity $v_A$) and the other as
the \emph{student} (denoted $B$, with trainable velocity $v_B=v_\theta$, parameters $\theta$).
The goal of distillation is to use the teacher's generative knowledge in $\Za$ to supervise the
student's generative process in $\Zb$.

\section{Latent Transport Map}
\label{sec:transport}

Because $v_A$ and $v_B$ are defined on \emph{different} latent spaces, every teacher signal
(clean samples, velocities, transition distributions) lives in $\Za$, whereas the student's
rollout, loss, and gradients all live in $\Zb$. Moving the supervision signal across requires a
bidirectional map between the latent spaces.

\begin{definition}[Latent transport map]
\label{def:transport}
A latent transport map is a pair of functions
\[
\Tmap:\Za\to\Zb,\qquad \Tinv:\Zb\to\Za,
\]
satisfying the mutual-inverse relations $\Tinv\circ\Tmap=\mathrm{id}_{\Za}$ and
$\Tmap\circ\Tinv=\mathrm{id}_{\Zb}$ (when the dimensions differ, $d_A\neq d_B$, invertibility
can only be required on a common subspace / range, i.e.\ as a generalized inverse).
\end{definition}

\begin{definition}[Semantic consistency / pixel anchoring]
\label{def:consistency}
We call $\Tmap$ \emph{pixel-consistent} if it preserves the decoded image semantics: for a
teacher latent $z\in\Za$,
\begin{equation}
\Dec_B\big(\Tmap(z)\big)\;\approx\;\Dec_A(z),
\qquad \forall z\in\Za.
\label{eq:pixel-consistency}
\end{equation}
That is, $z$ and $\Tmap(z)$ ``represent the same image''. This is the semantic-correctness
criterion for any reasonable transport.
\end{definition}

\begin{definition}[Transport consistency error]
\label{def:delta}
Define the round-trip pixel-space reconstruction error of the transport,
\begin{equation}
\delta_{\Tmap}
\;\coloneqq\;
\E_{z\sim p_A}\,\norm{\Dec_B\big(\Tmap(z)\big)-\Dec_A(z)},
\label{eq:delta}
\end{equation}
where $p_A$ is the teacher's latent distribution at the relevant time. $\delta_{\Tmap}=0$ means
exact pixel consistency \eqref{eq:pixel-consistency}. $\delta_{\Tmap}$ is the scalar that
quantifies the lossiness of a transport.
\end{definition}

\section{Pushforward and Pullback of the Velocity Field}
\label{sec:pushforward}

Distillation must transport not only clean samples but the \emph{entire velocity field}, so that
the teacher's probability flow in $\Za$, after being pushed through $\Tmap$, agrees with the
dynamics the student sees in $\Zb$.

\begin{problem}[Transport of the velocity field]
\label{prob:pf-existence}
Find $(\Tmap,\Tinv)$ such that the teacher's probability-flow ODE \eqref{eq:pf-ode} is
faithfully transported under the map $z_B(t)=\Tmap\big(z_A(t)\big)$.
\end{problem}

Formally, if $\Tmap$ is a $C^1$ diffeomorphism with Jacobian
$\Jac_{\Tmap}(z)\coloneqq\partial\Tmap/\partial z$, the chain rule gives
$\dot z_B=\Jac_{\Tmap}(z_A)\,\dot z_A$. Motivated by this, we \emph{define} the
\emph{pushforward velocity field} (the transported field expressed in $z_B$ coordinates) by
\begin{equation}
(\Tmap\push v_A)(z_B,t)
\;\coloneqq\;
\Jac_{\Tmap}\!\big(\Tinv(z_B)\big)\;
v_A\!\big(\Tinv(z_B),t\big),
\qquad z_B\in\Zb.
\label{eq:pushforward}
\end{equation}
The chain-rule identity above is the derivation that justifies this as the correct definition.
Symmetrically, mapping a state in the student space back into the teacher space requires the
inverse $\Tinv$ (pullback). Equation \eqref{eq:pushforward} makes explicit that, to evaluate the
teacher's ``opinion'' at a state visited by the student, one must simultaneously have:
\begin{enumerate}[leftmargin=2em,itemsep=2pt]
  \item \textbf{Invertibility}: access to $\Tinv(z_B)$, so the teacher can be queried at the
        student state $z_B$;
  \item \textbf{Differentiability}: the Jacobian $\Jac_{\Tmap}$ exists and is computable, so the
        teacher velocity can be correctly mapped into $\Zb$.
\end{enumerate}

\section{Dimension Mismatch and Conditional Pullback}
\label{sec:cond-pullback}

\textbf{Dimension (mis)match in practice.}
Empirically, the teacher's VAE has \emph{at least} as many channels as the student's VAE, i.e.
\[
d_A \ge d_B ,
\]
and strictly more in the common case. Two regimes arise:
\begin{itemize}[leftmargin=1.6em,itemsep=2pt]
  \item \emph{Strict mismatch} ($d_A>d_B$): the forward transport $\Tmap:\Za\to\Zb$ is a
        \emph{dimensionality reduction} (many-to-one), while the pullback $\Tinv:\Zb\to\Za$ is a
        \emph{dimensional expansion}; recovering a high-dimensional vector from a lower-dimensional
        one is \emph{underdetermined}, so the student $\to$ teacher transport is
        \textbf{necessarily lossy}.
  \item \emph{Equal dimension} ($d_A=d_B$): the dimension count alone no longer forces loss; if
        $\Tmap$ is full-rank (invertible) the single-valued pullback is exact, and the
        constructions below degenerate gracefully to the plain inverse of \S\ref{sec:transport}
        (made precise in Rem.~\ref{rem:equal-dim}).
\end{itemize}
This section treats the general case $d_A\ge d_B$, makes the strict-mismatch lossiness precise, and
formalizes the operation of ``conditioning on the student DiT's hidden states''.

\begin{definition}[Fiber and ideal pullback]
\label{def:fiber}
For $z_B\in\Zb$, the \emph{fiber} (preimage) of the forward transport is
\[
\mathcal{F}(z_B)\;\coloneqq\;\Tmap^{-1}(z_B)\;=\;\{\,z\in\Za:\ \Tmap(z)=z_B\,\}.
\]
When $\Tmap$ is a full-rank submersion, a generic fiber is a $(d_A-d_B)$-dimensional submanifold of
$\Za$ (an affine subspace when $\Tmap$ is affine); it has \emph{positive} dimension exactly when
$d_A>d_B$, and collapses to a \emph{single point} when $d_A=d_B$. Any \emph{single-valued} pullback
$\Tinv:\Zb\to\Za$ must select a \emph{unique} representative on each fiber, thereby discarding the
in-fiber coordinate---a vacuous choice when the fiber is a point ($d_A=d_B$), and a genuinely lossy
one when it is positive-dimensional ($d_A>d_B$). Moreover, let the teacher latent that is
pixel-consistent
(Def.~\ref{def:consistency}) with $z_B$ be the \emph{ideal pullback}
\begin{equation}
\bar z_A(z_B)\;\coloneqq\;\Enc_A\big(\Dec_B(z_B)\big)\ \in\ \Za,
\label{eq:ideal-lift}
\end{equation}
the semantically correct target that the pullback should reproduce.
\end{definition}

\begin{problem}[Irreducible loss of the underdetermined pullback]
\label{prob:underdetermined}
Suppose a single-valued pullback $\Tinv$ has range contained in some $d_B$-dimensional subset
$\mathcal{R}\subseteq\Za$ (e.g.\ a linear pullback $\Tinv(z_B)=A^{+}(z_B-b)$ lands in the column
space of $A^{+}$), and let $\Pi_{\mathcal R}$, $\Pi^{\perp}_{\mathcal R}$ be the corresponding
orthogonal projectors. Then for \emph{any} pullback that depends on $z_B$ alone,
\begin{equation}
\E_{z_B}\,\norm{\bar z_A(z_B)-\Tinv(z_B)}
\;\ge\;
\E_{z_B}\,\norm{\Pi^{\perp}_{\mathcal R}\,\bar z_A(z_B)},
\label{eq:irreducible}
\end{equation}
i.e.\ the energy of the ideal pullback lying in the \emph{teacher-exclusive subspace}
$\mathcal{R}^{\perp}$ (the extra $d_A-d_B$ dimensions) cannot be recovered from $z_B$ alone. This
is the precise meaning of the student $\to$ teacher transport being ``underdetermined and lossy''.
\end{problem}

\begin{remark}[The equal-dimension boundary $d_A=d_B$]
\label{rem:equal-dim}
When $d_A=d_B$ and the chosen $\Tmap$ is invertible, the recoverable subspace is the whole space,
$\mathcal{R}=\Za$, hence $\Pi^{\perp}_{\mathcal R}=0$ and the right-hand side of
\eqref{eq:irreducible} vanishes: a single-valued pullback can be exact and there is \emph{no}
dimension-induced loss. The fiber is a point, so $\bar z_A(z_B)=\Tinv(z_B)$, and the conditional
pullback introduced below carries no information-theoretic gain
($I(\bar z_A;h_\theta\mid z_B)=0$, since $\bar z_A$ is then a deterministic function of $z_B$);
requirement (P4) thus collapses to (P2). This concerns only the \emph{dimension-induced}
underdetermination: a residual approximation error may still persist if the chosen transport class
(e.g.\ affine) cannot represent $\bar z_A$, but that error is measured by $\delta_{\Tmap}$
(Def.~\ref{def:delta}), not by \eqref{eq:irreducible}. In short, $d_A=d_B$ is the degenerate
boundary of this section at which everything reduces to \S\ref{sec:transport}--\S\ref{sec:pushforward}.
\end{remark}

\textbf{Conditional pullback: injecting information via the student DiT's hidden states.}
While computing $v_\theta(z_B,t,c)$ (with $c$ the external conditioning, e.g.\ the text prompt),
the student transformer produces intermediate hidden states
\[
h_\theta(z_B,t,c)\ \in\ \mathcal{H},
\]
whose dimension is typically far larger than $d_B$. Augmenting the pullback's domain from $\Zb$
to $\Zb\times\mathcal{H}$ yields the \emph{conditional pullback}
\begin{equation}
\Tinv_{c}:\ \Zb\times\mathcal{H}\to\Za,
\qquad
\hat z_A\;\coloneqq\;\Tinv_{c}\big(z_B,\ h_\theta(z_B,t,c)\big).
\label{eq:cond-pullback}
\end{equation}

\begin{remark}[Essence: turning an underdetermined inversion into conditional inference]
\label{rem:essence}
The essence of this operation is to use the student's own internal representation as
\emph{side information} to estimate the missing coordinate in the teacher-exclusive subspace
$\mathcal{R}^{\perp}$, thereby rewriting an ``ill-posed dimensional-expansion inversion'' as a
``conditional inference / conditional completion'' problem. Its attainable accuracy is governed
by the \emph{conditional entropy}:
\begin{equation}
H\big(\bar z_A \,\big|\, z_B,\ h_\theta\big)
\;\le\;
H\big(\bar z_A \,\big|\, z_B\big),
\qquad
\text{gain}\;=\;I\big(\bar z_A;\ h_\theta \,\big|\, z_B\big)\;\ge\;0 .
\label{eq:cond-entropy}
\end{equation}
Under squared error, the optimal estimator is the conditional expectation
$\hat z_A^{\star}=\E[\bar z_A\mid z_B,h_\theta]$, whose irreducible residual is the conditional
variance $\E\,\mathrm{Var}\big(\Pi^{\perp}_{\mathcal R}\bar z_A\mid z_B,h_\theta\big)$.
\end{remark}

\begin{remark}[A necessary caveat: hidden states do not create information from nothing]
\label{rem:determinism}
Since $h_\theta(z_B,t,c)$ is a \emph{deterministic} function of $(z_B,t,c)$, the genuine
information gain from conditioning can \emph{only} come from inputs beyond $z_B$, namely the time
$t$ and the external conditioning $c$:
\[
I\big(\bar z_A;\ h_\theta \,\big|\, z_B\big)
\;\le\;
I\big(\bar z_A;\ (t,c)\,\big|\, z_B\big),
\]
which is strictly positive iff the teacher-exclusive coordinate
$\Pi^{\perp}_{\mathcal R}\bar z_A$ is correlated with $(t,c)$. The role of the hidden states is
therefore twofold: (i) they are the natural carrier of the $(z_B,t,c)$ context (in particular the
prompt $c$, which is \emph{not} a function of $z_B$); and (ii) they are a precomputed, rich
nonlinear feature that makes the conditional completion learnable in practice. They \emph{cannot}
recover the part of the information that was \emph{physically deleted} by $\Tmap$ in the forward
direction and is uncorrelated with $(t,c)$.
\end{remark}

\begin{remark}[Compatibility with the pushforward]
\label{rem:compat}
Substituting the conditional pullback into the pushforward \eqref{eq:pushforward}:
\[
(\Tmap\push v_A)(z_B,t)
\;=\;\Jac_{\Tmap}(\hat z_A)\,v_A(\hat z_A,t),
\qquad \hat z_A=\Tinv_{c}\big(z_B,h_\theta(z_B,t,c)\big).
\]
For an affine $\Tmap=Az+b$, the Jacobian $\Jac_{\Tmap}\equiv A$ is constant, so conditioning
\emph{does not alter} the velocity-transport operator $A$; it only \emph{relocates} the teacher
query point $\hat z_A$ to a more correct position within the fiber $\mathcal{F}(z_B)$. That is, it
sharpens \emph{which} teacher opinion is queried without breaking the linear structure of the
transport.
\end{remark}

\section{Formal Requirements of On-Policy Distillation}
\label{sec:onpolicy}

The defining feature of on-policy distillation is that the supervision signal must be evaluated on
\emph{the state distribution that the student itself rolls out}, rather than on states presampled
by the teacher.

\paragraph{The student's on-policy states.}
The student rolls out a trajectory $\{z_B(t)\}_{t}\subset\Zb$ using its own velocity $v_\theta$,
inducing a distribution $\rho_\theta$ of visited states (it depends on the current parameters
$\theta$, hence ``on-policy'').

\paragraph{The supervision signal.}
For each visited state $z_B\sim\rho_\theta$, distillation requires the teacher's ``opinion'' on
how to denoise at that point. But the teacher velocity $v_A$ is defined only on $\Za$, so its
legitimate value in $\Zb$ is exactly the pushforward field $(\Tmap\push v_A)(z_B,t)$ of
\eqref{eq:pushforward}, which depends on $\Tinv$ (pulling $z_B$ back to $\Za$ to query the
teacher) and $\Jac_{\Tmap}$ (mapping the result back into $\Zb$).

\paragraph{Distillation objective (general form).}
On-policy distillation minimizes the discrepancy between the student velocity and the transported
teacher velocity over the student's visited distribution:
\begin{equation}
\min_{\theta}\;
\E_{t}\,\E_{z_B\sim\rho_\theta}\;
\Big[\,w(t)\,
\mathcal{M}\!\Big(v_\theta(z_B,t),\;(\Tmap\push v_A)(z_B,t)\Big)\Big],
\label{eq:distill-objective}
\end{equation}
where $w(t)$ is a time weight and $\mathcal{M}(\cdot,\cdot)$ is a discrepancy measure (e.g.\ a
quadratic regression error, or a divergence between the transition distributions induced by the
dynamics); its concrete form is outside the scope of this note. The objective is well defined
precisely when $(\Tmap\push v_A)$ can be evaluated on the support of $\rho_\theta$, i.e.\ when
$\Tinv$ and $\Jac_{\Tmap}$ exist and are computable.

\section{Formal Problem Statement}
\label{sec:problem-statement}

Combining \S\ref{sec:setup}--\S\ref{sec:onpolicy}, cross-latent on-policy distillation reduces to
the following transport-map design problem.

\begin{problem}[Cross-latent transport design]
\label{prob:main}
Given two heterogeneous latent spaces $\Za,\Zb$, their decoders $\Dec_A,\Dec_B$, the teacher
velocity $v_A$, and the student velocity $v_\theta$, find a pair of maps $(\Tmap,\Tinv)$,
$\Tmap:\Za\to\Zb$ and $\Tinv:\Zb\to\Za$, such that:
\begin{enumerate}[leftmargin=2.4em,label=\textup{(P\arabic*)},itemsep=3pt]
  \item \textbf{Semantic consistency}: $\Tmap$ is pixel-consistent, i.e.\ the transport
        consistency error $\delta_{\Tmap}$ (Def.~\ref{def:delta}) is as small as possible,
        $\Dec_B(\Tmap(z))\approx\Dec_A(z)$;
  \item \textbf{Invertibility}: $\Tinv$ satisfies the mutual-inverse relations of
        Def.~\ref{def:transport}, so student states can be pulled back to $\Za$ to query the
        teacher;
  \item \textbf{Differentiable pushforward}: $\Tmap$ is differentiable and its Jacobian
        $\Jac_{\Tmap}$ is available, so the pushforward velocity field
        \eqref{eq:pushforward} $(\Tmap\push v_A)$ can be evaluated at the student-visited states;
  \item \textbf{Conditional invertibility under mismatch} (active only for $d_A>d_B$): with
        $d_A\ge d_B$, the loss of a single-valued pullback $\Tinv(z_B)$ is lower-bounded by
        \eqref{eq:irreducible}. When $d_A>d_B$ this bound is generally positive (underdetermined
        expansion, necessarily lossy), and one should use the conditional pullback
        $\Tinv_{c}\big(z_B,h_\theta(z_B,t,c)\big)$ (\S\ref{sec:cond-pullback}), using the student
        DiT's hidden states as side information to complete the teacher-exclusive subspace
        $\mathcal{R}^{\perp}$, so the pullback error approaches the conditional-variance lower
        bound \eqref{eq:cond-entropy}. When $d_A=d_B$ the bound vanishes and this requirement
        reduces to (P2) (Rem.~\ref{rem:equal-dim}).
\end{enumerate}
When (P1)--(P4) hold, the supervision signal of the on-policy distillation objective
\eqref{eq:distill-objective} is type-legal and semantically correct, and the distillation problem
becomes well posed.
\end{problem}

\begin{remark}[Hierarchy of requirements]
\label{rem:layering}
Different strengths of distillation impose different requirements on $(\Tmap,\Tinv)$: performing
only \emph{sample transport} $z_B^{0}\coloneqq\Tmap(z_A^{0})$ on clean samples $z_A^{0}$ requires only
(P1), whereas performing \emph{velocity-field transport} on the student's on-policy trajectory
additionally requires (P2)--(P3), and, under \emph{strict} dimension mismatch $d_A>d_B$, (P4)
(at $d_A=d_B$, (P4) reduces to (P2)). These levels constitute the several strength variants of
Problem~\ref{prob:main}.
\end{remark}

\begin{remark}[Scope statement]
This note gives only the problem formalization (setup, required properties, problem statement). It
does not discuss any concrete construction, parameterization, initialization, or solution algorithm
for $(\Tmap,\Tinv)$, nor does it compare or select among candidate methods.
\end{remark}

\begin{remark}[External concrete $\Tinv$: Cross-Space Bridge]
\label{rem:bridge}
\href{https://arxiv.org/abs/2606.32020}{Cross-Space Distillation (arXiv:2606.32020)} learns a
Bridge $\mathcal{B}_\phi:\Zb\to\Za$ (Student-decoder spatial prior $+$ projector) so that
distribution-based one-step distillation (VSD/GAN) can run under Cross-Resolution and Cross-VAE
mismatch. In our terms this is a concrete $\Tinv$ satisfying a soft version of (P1)--(P2) for
\emph{sample} transport into Teacher space; it does not by itself supply the Jacobian pushforward
(P3) needed for Student-space velocity matching. See
\texttt{docs/xopd/cross\_vae/cross\_space\_distillation\_bridge.tex} for the mapping onto Routes~A--C and
the XOPD L0/L1 hierarchy (Rem.~\ref{rem:layering}).
\end{remark}

\end{document}
